\section{Jahwe Zebaoth}
\begin{block}[Namen Gotte in der Bibel]
        Letzten Sonntag haben wir die drei Hauptnamen, Jehowa oder Jahweh, Elohim und Adonai plus ein paar Eigenschaften von Gott angeschaut. Oft werden für Gott in der Bibel Namen so zusammengesetzt, dass sie seine Eigenschaften ausdrücken. Ein paar dieser Namen haben wir uns letzten Sonntag angeschaut. Eigentlich wollte ich Heute etwas über die Dreieinigkeit Gottes sagen, aber ich fand wir sollten uns doch noch ein bisschen mit den Namen beschäftigen. Schon bei dem ersten Namen kamen mir viele Schätze entgegen. Weiss nicht wieviele der Namen wir anschauen werden, aber heute wollen wir uns als erstes den Namen \betonung[b,p]{Jahweh Zebaoth}  widmen.

        Dieser Name kommt in der Bibel in ca. 274 Versen vor. Der Name kommt zum erstenmal in 1. Samuel vor. Danach vorwiegend bei den Propheten. So schreibt Jeremia über 80 mal \enquote{so spricht der HERR Zebaoth}, Jesaja und Sachaja zusammen reden Gott über 100 mal mit diesem Titel an. In den Psalmen wird der Name 15 mal erwähnt. 

        Dieser Name Gottes kann man als \enquote{Herr der Heerscharen} übersetzen. Das heisst Gott ist Herr über mehr als nur über uns Menschen. Wir wissen aus der Bibel, dass Gott mehr geschaffen hat, als nur uns Menschen und die Tiere auf dieser Erde. In der Bibel kommen diverse Himmelswesen vor, so \zB{} die Engel und die Cherubim. In Hesekiel werden Wesen mit 4 Gesichtern beschrieben oder Wesen die aussehen wie Wagenräder voller Augen. Auch der Satan wurde von Gott geschaffen. Gott ist Herr über das alles und über uns Menschen. Gott ist Herrscher über das ganze Universum und über die Himmel. Dies soll dieser Name ausdrücken. Der Name betont die Eigenschaften: \betonung[b,p]{Macht, Souveränität, Schutz und König aus.}
    \end{block}
    \begin{block}[Jahweh Zebaoth]
        Wie erkläre ich euch jetzt am besten den Namen? Am einfachsten lasse ich Gottes Wort sprechen.
        So wollen wir uns den Psalm 24 zusammen anschauen, der erklärt am besten  was dieser Name bedeutet. Lesen wir zusammen diesen Psalm 24.
        
        Wie Vers 1 uns sagt, wurde dieser Psalm von David geschrieben. Es ist ein Lied, das uns erklärt wie Gross unser Gott ist. In den ersten drei Versen zeigt uns David, dass dieser Gott unerreichbar für uns Menschen ist. Nur ausgewählten Personen durften auf den Berg Sinai, alle anderen -- Mensch und Tier -- mussten sterben, wenn sie nur den Berg berührten. Auch das Allerheiligste in der Stiftshütte oder später im Tempel durfte nur von ausgewählten Priester betreten werden, sonst musste man sterben. 
        
        In den Versen 4 - 5 sagt uns das Lied, wie man zum Herrn kommen kann. Man muss Unschuldig sein, keine Lügen in sich haben. Und sind wir das? Sind wir ohne Lug und Trug? Können wir als sündige Menschen einfach so vor den Herrn treten? 

        In Vers 6 sagt uns David, dass nur das Volk Israel nach diesem Gott sucht. Er sagt nicht, dass die Juden direkt zum Herrn dürfen, sondern das dieses als einziges Volk den Herrn sucht. In den restlichen Versen erklärt David nochmals wie Mächtig und wie Herrlich dieser Gott ist. In Vers 10 dann der Höhepunkt:
        \begin{bibelbox}{SCHL}{Ps}{24:10}
            Wer ist denn dieser König der Herrlichkeit? Jahweh Zebaoth (Der Herr der Heerscharen), \betonung[b,p]{er} ist der König der Herrlichkeit.
        \end{bibelbox}
        Dieser Psalm zeigt uns die absolute Macht, Souveränität und Heiligkeit Gottes. Kein Mensch ist gut genug für diesen Gott. Aus diesem Grund können wir nicht mit Werken unser Heil verdienen. Weil aber Gott gerecht ist, hat er uns einen Ausweg aus diesem Dilemma geschafft. Er hat uns seinen Sohn als Sühneopfer geschenkt. Darum schreibt Paulus den Römer:
        \begin{bibelbox}{SCHL}{Rom}{10:9-11}
        Denn wenn du mit deinem Mund bekennst und in deinem Herzen glaubst, dass Gott ihn aus den Toten auferweckt hat, so wirst du gerettet. Denn mit dem Herzen glaubt man, um gerecht zu werden, und mit dem Mund bekennt man, um gerettet zu werden; denn die Schrift spricht: \enquote{Jeder, der an ihn glaubt, wird nicht zuschanden werden!}
        \end{bibelbox}
        Der Vater hat uns seinen eigenen Sohn geopfert. Nur wenn wir diese Opfer Jesus Christus von Herzen annehmen, werden wir vor Jahweh Zebaoth bestehen können. 

        Am Ende der Zeit wird dieser Herr Heerscharen diese Welt regieren in Friede und Gerechtigkeit. Freuen wir uns auf seine Wiederkunft.
    \end{block}